% ===========================================================================
\chapter{Further Contributions}
\label{chap:further}
\margintoc
% ===========================================================================

Four further publications carry the thesis author's name and are not
reproduced in this thesis. Each is summarised here in a page or two, with the
author's own role stated first, because the contribution to each was partial
or because the work sits beside the argument instead of carrying it. The survey
of the previous chapter is summarised for the same reason and on the same
terms. The published papers stand on their own and are cited in full. This
follows the University of Zurich and Institute of Neuroinformatics expectation
that a cumulative thesis separates the candidate's own intellectual
contribution from collaborative work, and it keeps the reproduced chapters to
the four papers the argument rests on.

Two of the four measure what current models do, and they sharpen the gaps the
previous chapter names. One supplies structure from outside the data and measures
what it buys, and one holds its knowledge outside the weights altogether.

% ---------------------------------------------------------------------------
\section{Quantisation costs text quality that the standard metrics miss}
\label{sec:llmbench}
% ---------------------------------------------------------------------------

\begin{publicationbasis}{Published}
Lukas Tuggener, Pascal Sager, Yassine Taoudi-Benchekroun, Benjamin F. Grewe
and Thilo Stadelmann. So you want your private LLM at home? A survey and
benchmark of methods for efficient GPTs. Proceedings of the 11th IEEE Swiss
Conference on Data Science, 2024, pages $205$ to
$212$~\cite{tuggener_you_2024}.
\end{publicationbasis}

The thesis author ran the evaluation of the language models. The survey, the
benchmark design and the write up are the work of the co-authors, and Lukas
Tuggener led the paper.

The paper asks what it costs to run a capable language model on hardware a
small group can afford. It reviews the methods that reduce the computational
footprint of a large language model and then benchmarks seven of them, drawn
from the two areas where most progress had been made, model quantisation and
low rank adapters. The benchmark reports resource consumption such as memory
on the graphics card, a quantitative performance evaluation, and a qualitative
study in which eight human raters judged the generated text.

Two results matter here. Quantisation reduces the memory a model needs by a
large factor, which is the saving it is sold on. It also costs text quality
that the human raters noticed and that the standard automatic metrics missed,
where the advertising of the day promised a saving without a cost. Low rank
adapters allow a model to be fine tuned with moderate compute and avoid that
cost.

The second result is the one the thesis uses. A method was accepted on a score
that was not sensitive to the property the score was taken to stand for, and
the gap only became visible when a different readout, human judgement, was
applied to the same models. \Cref{chap:bounded} analyses the same failure in
a setting where the readout admits a correction. The two are cases of one
mistake, with a different mechanism behind each.

% ---------------------------------------------------------------------------
\section{Architectures built for composition still fail on unseen rule combinations}
\label{sec:cogitao}
% ---------------------------------------------------------------------------

\begin{publicationbasis}{Published}
Yassine Taoudi-Benchekroun and colleagues. COGITAO: a visual reasoning
framework to study compositionality and generalization. 2nd ICML Workshop on
Compositional Learning, 2026~\cite{taoudibenchekroun_cogitao_2026}.
\end{publicationbasis}

The thesis author evaluated one model family, the masked latent diffusion
model, on the released benchmark. The framework, the data generator, the task
design and the remaining baselines are the work of the co-authors.

COGITAO is a data generator for studying whether a model can apply learned
concepts in combinations it has not seen. It builds rule based tasks on grids,
in the spirit of the ARC-AGI setting, by applying transformations to objects.
It supports $28$ transformations that can be composed with one another, at a
composition depth the experimenter sets, with control over the grid and the
properties of the objects on it. The number of distinct task rules this allows
runs into the millions, and any number of samples can be drawn per rule, so
the in domain and out of domain conditions separate by construction with no
need for held out sampling.

The released baselines include architectures chosen because their inductive
biases suit composition, among them a diffusion based Transformer and
recurrent Transformers with adaptive computation time. Every one of them
reaches strong accuracy on the compositions it was trained on, and every one
fails on compositions held out from training. The failure is consistent across
the architectures tested.

For this thesis the result is a measurement of the first gap named in
\cref{chap:acu}. A model that scores well in domain can still fail
consistently on a recombination of the same components, so in domain accuracy does not
report whether the components were learned as components. The benchmark does
not say what in training produces either outcome, and the thesis does not use
it to argue that it does.

% ---------------------------------------------------------------------------
\section{Tissue ranges from the medical literature buy accuracy and data efficiency}
\label{sec:hounsfield}
% ---------------------------------------------------------------------------

\begin{publicationbasis}{Published}
Benjamin Meyer, Pascal J. Sager and colleagues. Hounsfield unit ranges as
inductive bias for intra clinical learning of data efficient CT segmentation
models~\cite{meyer_hounsfield_2025}.
\end{publicationbasis}

The thesis author is a shared first author and contributed mainly the framing
of the work. The implementation and the experiments are the work of the
co-authors, and Benjamin Meyer led them.

Computed tomography measures radiodensity in Hounsfield units, and different
tissue types occupy different and largely known ranges of that scale. The
paper turns that knowledge into an input. For each tissue type of interest it
takes the range from the medical literature, builds a binary mask marking the
voxels inside the corresponding range, and concatenates those masks with the
original scan as extra input channels. The segmentation model is otherwise
unchanged.

The comparison is against the same model without the masks, trained on the
same data with the same budget. Adding the masks improves segmentation
accuracy, and it improves data efficiency, so the model reaches the baseline's
accuracy with less annotated training data. Annotation is the scarce resource
in a clinical setting, so the second result is the useful one.

The work belongs in this thesis as a case where structure was supplied
deliberately from outside the data and the gain was measured against the one
comparison that settles it, the identical model without that structure.

% ---------------------------------------------------------------------------
\section{Knowledge held in an index at inference time needs no training}
\label{sec:retrieval}
% ---------------------------------------------------------------------------

\begin{publicationbasis}{Published}
Pascal J. Sager, Ashwini Kamaraj, Benjamin F. Grewe and Thilo Stadelmann.
Deep Retrieval at CheckThat! 2025: identifying scientific papers from implicit
social media mentions via hybrid retrieval and re-ranking. CLEF 2025 Working
Notes~\cite{sager_deep_2025}.
\end{publicationbasis}

The thesis author is the first author and designed and built the retrieval
pipeline. Ashwini Kamaraj contributed to the experiments, and Benjamin F.
Grewe and Thilo Stadelmann supervised.

The task is to find the scientific paper a social media post refers to when
the post names no paper. The language of the post and the language of the
paper differ in register and in vocabulary, so lexical matching alone is
brittle and semantic matching alone is imprecise. The pipeline uses both.
Keyword matching with BM25 returns the $30$ best candidates, a dense search
over a vector store using a fine tuned retriever returns $100$, the two
candidate sets are merged, and a cross encoder built on a language model
rescores what is left.

The pipeline reaches a mean reciprocal rank at five of $76.46\%$ on the
development set, which placed first on that leaderboard, and $66.43\%$ on the
hidden test set, which placed third of $31$ teams and within two percentage
points of the winning system. It runs open source models locally and uses no
external training data.

This work sits beside the thesis argument because of the form its solution
takes. The knowledge the system needs is held in an index and fetched while
the system runs, with none of it written into the weights. That is one way
around the question the rest of the thesis asks, namely what training writes
into a network and where.
